\section{Evaluation of the benchmark}\label{sec:analysis}
The results in this section estimate the explicit sequence $L_n$ in
\eqref{eq:normalization}. Theorem~\ref{thm:main} transfers them to the subgroup count with an
exponentially smaller relative error.

\subsection{The binary Gaussian sum}
Define
\begin{equation}\label{eq:theta-constants}
 \kappa_0=\phi^{-1}\sum_{j\in\Z}2^{-j^2},\qquad
 \kappa_1=\phi^{-1}\sum_{j\in\Z}2^{-j(j-1)}.
\end{equation}
Both series converge absolutely.
\begin{lemma}\label{lem:Gaussian-asymptotic}
As $R\to\infty$,
\[
 G_R=\kappa_{R\bmod2}\,2^{\lfloor R^2/4\rfloor}
            \bigl(1+O(R2^{-R/2})\bigr).
\]
\end{lemma}
\begin{proof}
The product formula reads
\[
 \gauss Rk2=2^{k(R-k)}
 \frac{\prod_{i=R-k+1}^R(1-2^{-i})}
      {\prod_{i=1}^k(1-2^{-i})}.
\]
For central $k$ the fraction is
$\phi^{-1}(1+O(2^{-k}+2^{-(R-k)}))$.
When $R=2m$, put $k=m+j$, giving exponent $m^2-j^2$.
When $R=2m+1$, put $k=m+j$, giving exponent
$m(m+1)-j(j-1)$. On $|j|\leq R/4$ the summable Gaussian weights
times $2^{|j|}$ give error $O(2^{-R/2})$ up to an absolute parity
constant. The remaining tails are $2^{-\Omega(R^2)}$ by
\eqref{eq:gaussian-bounds}. The stated, slightly weaker, error follows.
\end{proof}

\subsection{One positive saddle}
Let $\rho=\rho_R>0$ be the unique solution of
\begin{equation}\label{eq:saddle}
 R=\rho P'(\rho)=\frac\rho2+\frac{\rho^2}3+\frac{\rho^4}{96},
 \qquad b(\rho)=\frac\rho2+\frac{2\rho^2}3+\frac{\rho^4}{24}.
\end{equation}
Existence and uniqueness follow from strict monotonicity on the positive
axis. In particular $\rho\sim(96R)^{1/4}$ and $b(\rho)\sim4R$.

\begin{theorem}[Saddle formula]\label{thm:saddle}
For $n=2R+\eps$,
\begin{equation}\label{eq:saddle-formula}
 L_n=n!G_R\,
 \frac{(1+\rho/6)^\eps e^{P(\rho)}}
      {\rho^R\sqrt{2\pi b(\rho)}}
       \left(1+O(R^{-1})\right).
\end{equation}
In particular $c_{R-1}/c_R=\rho(1+O(R^{-1}))$.
\end{theorem}
\begin{proof}
Put $\lambda_1=\rho/2$, $\lambda_2=\rho^2/6$ and
$\lambda_4=\rho^4/384$, so that
$R=\sum_k k\lambda_k$ and $b=\sum_k k^2\lambda_k$ satisfy
$R\leq b\leq4R$. Here and below $k\in\{1,2,4\}$. Define
\[
 A(\theta)=\sum_k\lambda_k(1-\cos k\theta),\qquad
 B(\theta)=\sum_k\lambda_k(\sin k\theta-k\theta).
\]
Then $P(\rho e^{i\theta})-P(\rho)-iR\theta=-A+iB$.
For any $d\in[0,1]$, set
\[
 F_d(\theta)=(1-d)\cos B(\theta)+d\cos(\theta+B(\theta)),\qquad
 H_d(\theta)=e^{-A(\theta)}F_d(\theta).
\]
Taking real parts in Cauchy's formula on $y=\rho e^{i\theta}$ gives
\[
 \frac{\rho^R[y^R](1+y/6)^\eps e^{P(y)}}
      {(1+\rho/6)^\eps e^{P(\rho)}}
 =\frac1{2\pi}\int_{-\pi}^{\pi}H_d(\theta)\,d\theta,
 \qquad d=\begin{cases}0,&\eps=0,\\ \rho/(6+\rho),&\eps=1.\end{cases}
\]
We estimate this integral uniformly for all $d\in[0,1]$.

Use the fixed central interval $|\theta|\leq\alpha=\pi/4$ and write
$Q=b\theta^2/2$. The inequalities
$u^2/8\leq1-\cos u\leq u^2/2$ for $|u|\leq\pi$, applied to each
$k\theta$, give
\[
 R\theta^2/8\leq A\leq Q,\qquad
 0\leq Q-A=O(R\theta^4),\qquad |B|=O(R|\theta|^3).
\]
The last two bounds follow from the fourth-order cosine and third-order
sine remainders; their constants are absolute. Since $F_d$ is a convex
combination of cosines,
\[
 |F_d|\leq1,\qquad
 0\leq1-F_d\leq B^2/2+|\theta B|+\theta^2/2.
\]
Together with
$0\leq e^{-A}-e^{-Q}\leq e^{-A}(Q-A)$, this yields
\[
 |H_d(\theta)-e^{-b\theta^2/2}|
 \leq C e^{-R\theta^2/8}
       \bigl(\theta^2+R\theta^4+R^2\theta^6\bigr).
\]
Integrating this bound gives $O(R^{-3/2})$, by the change of variable
$u=\sqrt R\,\theta$. This argument requires no smallness assumption on $B$.

On the entire complement $\alpha\leq|\theta|\leq\pi$, the linear term
alone gives
\[
 A(\theta)\geq\frac\rho2(1-\cos\alpha)\geq\frac\rho{16},\qquad
 |H_d(\theta)|\leq e^{-\rho/16}.
\]
Because $\rho\asymp R^{1/4}$, these tails are smaller than every fixed
inverse power of $R$. The Gaussian tails outside $[-\alpha,\alpha]$
are $O(e^{-cR})$ for an absolute $c>0$. Thus, uniformly in $d$,
\[
 \int_{-\pi}^{\pi}H_d(\theta)\,d\theta
 =\sqrt{\frac{2\pi}{b}}+O(R^{-3/2})
 =\sqrt{\frac{2\pi}{b}}\bigl(1+O(R^{-1})\bigr),
\]
which proves \eqref{eq:saddle-formula}.
Finally, Cauchy's formula with amplitude $y/\rho$ gives $H_1$, at the
same saddle $\rho_R$. Comparing it with $H_0$ gives
\[
 \frac{c_{R-1}}{\rho c_R}
 =\frac{\int_{-\pi}^{\pi}H_1(\theta)\,d\theta}
        {\int_{-\pi}^{\pi}H_0(\theta)\,d\theta}
 =1+O(R^{-1}).\qedhere
\]
\end{proof}

\subsection{An elementary four-periodic expression}
Set
\begin{equation}\label{eq:four-constants}
 \begin{aligned}
 C_0&=e^{-4/3}\kappa_0/2^{1/2},&
 C_2&=e^{-4/3}\kappa_1/2^{3/4},\\
 C_1&=\frac{e^{-4/3}\kappa_0\,48^{3/8}}{6\cdot2^{7/16}},&
 C_3&=\frac{e^{-4/3}\kappa_1\,48^{3/8}}{6\cdot2^{11/16}}.
 \end{aligned}
\end{equation}
For $\eps=n\bmod2$ define
\begin{equation}\label{eq:elementary-M}
 M_n=C_{n\bmod4}\,n^{3\eps/8}\,2^{n^2/16}
 \left(\frac{n^7}{48\cdot2^\eps e^7}\right)^{n/8}
 \exp\left(2\sqrt{n/3}+\frac{(48n)^{1/4}}2\right).
\end{equation}

\begin{theorem}[Elementary expansion]\label{thm:elementary}
The exact benchmark satisfies
\begin{equation}\label{eq:elementary-error}
 \frac{L_n}{M_n}=1+\frac{6\eps-4}{(48n)^{1/4}}+O(n^{-1/2}).
\end{equation}
By Theorem~\ref{thm:main}, the same formula
holds with $s_n$ in place of $L_n$, and
\eqref{eq:saddle-formula} holds with $s_n$ on the left.
\end{theorem}
\begin{proof}
Put $x=(96R)^{1/4}$ and $v=x(x-\rho)$. Only the first displacement
of the saddle is needed:
\begin{equation}\label{eq:rho-expansion}
 \rho=x-8x^{-1}+O(x^{-2}),\qquad v=8+O(x^{-1}).
\end{equation}
Indeed, the equation is $\rho^4+32\rho^2+48\rho=x^4$.
For large $x$ its positive root lies in $[x/2,x]$; substituting
$x-8/x$ leaves residual $O(x)$, and the derivative throughout the
intervening interval is of order $x^3$. The mean value theorem proves
\eqref{eq:rho-expansion}.

Expand $\log_e\rho=\log_ex+\log_e(1-v/x^2)$ to quadratic order.
Since $v$ is bounded, the cubic remainder, multiplied by $R=x^4/96$,
is $O(x^{-2})$. Direct substitution gives
\begin{align*}
 P(\rho)-R\log_e\rho
 &=\frac R4-R\log_ex+\frac{x^2}{6}+\frac x2
       +\frac{v^2}{48}-\frac v3-\frac{v}{2x}+O(x^{-2})\\
 &=\frac R4-R\log_ex+\frac{x^2}{6}+\frac x2
       -\frac43-\frac4x+O(x^{-2}).
\end{align*}
The cancellation is explicit: $v^2/48-v/3=-4/3+(v-8)^2/48$.
Also $b(\rho)/(4R)=1+O(x^{-2})$.

Define
\[
 T_R=\frac{1}{\sqrt{8\pi R}}\exp\left(
 -\frac R4\log_eR+\frac{1-\log_e96}{4}R
 +\frac{2\sqrt6}{3}\sqrt R+\frac x2-\frac43\right).
\]
The preceding estimates and the saddle formula yield
\begin{equation}\label{eq:log-coefficient-expansion}
 \log_e\frac{c_R}{T_R}=-\frac4x+O(x^{-2}).
\end{equation}
The relative saddle error is $O(R^{-1})=O(x^{-4})$ and is absorbed
in this remainder. For the odd amplitude, the same first displacement gives
\begin{equation}\label{eq:odd-amplitude-expansion}
 \log_e\frac{1+\rho/6}{x/6}=\frac6x+O(x^{-2}).
\end{equation}

To make the parity conversion explicit, set $y=(48n)^{1/4}$ and
\begin{align*}
 \Lambda_n={}&(n+\tfrac12)\log_en-n+\tfrac12\log_e(2\pi)
   +\lfloor R^2/4\rfloor\log_e2+\log_e\kappa_{R\bmod2}\\
 &+\log_eT_R+\eps\log_e(x/6)+(6\eps-4)/x.
\end{align*}
Stirling's formula, Lemma~\ref{lem:Gaussian-asymptotic} and the two
expansions above give $\log_e L_n=\Lambda_n+O(n^{-1/2})$.
Using $R=(n-\eps)/2$ and the exact constants in
\eqref{eq:four-constants}, elementary algebra gives
\begin{align*}
 \Lambda_n-\log_eM_n-\frac{6\eps-4}{y}
 ={}&\frac{n-3\eps+4}{8}\log_e\frac{n}{n-\eps}-\frac\eps8\\
 &+\frac{x^2-y^2}{6}+\frac{x-y}{2}
       +(6\eps-4)\left(\frac1x-\frac1y\right).
\end{align*}
In even degree this is identically zero. In odd degree the first line
is $O(n^{-1})$, while $y^4-x^4=48$ bounds the remaining three
differences by $O(n^{-1/2})$, $O(n^{-3/4})$ and $O(n^{-5/4})$,
respectively. Thus
\[
 \log_e(L_n/M_n)=(6\eps-4)(48n)^{-1/4}+O(n^{-1/2}).
\]
Exponentiation proves \eqref{eq:elementary-error}. Finally the
exponential relative error of Theorem~\ref{thm:main}
is smaller than either polynomial error asserted here.
\end{proof}

\begin{remark}
The relative errors refer to three different approximations. The
complete counting theorem gives exponential error against the
exact coefficient benchmark $L_n$. Replacing the coefficient by its
single-saddle expression gives $O(n^{-1})$. Replacing the saddle by
the elementary expression $M_n$ gives an explicit term of order
$n^{-1/4}$ and a remainder $O(n^{-1/2})$.
\end{remark}
